\section{Return to the four-ball}\label{sec:rtc}

We use the torus cover of \S\ref{sec:fer} and unit representatives in $\C^3\times\C$. Write $\mathcal F(Z)=Z_0^3+Z_1^3+Z_2^3$, fix $0<h<d_n<1/3$, and put $\mathcal N_h=\{(Z,U):|Z|=1,\ \mathcal F(Z)=0,\ \pi|U|^2\le h\}$. The cover map $\mathsf N(u,v,w)=(Z(u,v),w)$ takes $\{p\le h\}$ into $\mathcal N_h$. The three marked base points are the zeros of $Z_0$ on the torus; on the cover, the zero set is their $\Z^2$-periodic inverse image (Lemma~\ref{lem:fer-bpoints}). Pairs are gauge equivalent if $(Z',U')=(cZ,c^3U)$ with $|c|=1$.

\subsection{Conventions and the ambient form}\label{sec:rtc-conventions}

For $Z\ne0$ put $\eta_H=\operatorname{Im}(Z^*dZ)/(2\pi|Z|^2)$ and $\omega_{\rm FS}=d\eta_H=dd^c\log|Z|^2$. For unit $Z$ put $H_Z=\{\xi:Z^*\xi=0\}$. On $H_Z$ one has $\omega_{\rm FS}(\xi_1,\xi_2)=\pi^{-1}\operatorname{Im}(\xi_1^*\xi_2)$, because the terms of $d\eta_H$ containing $d|Z|^2=2\operatorname{Re}(Z^*\,\cdot\,)$ vanish there. All forms below are gauge invariant and annihilate the gauge generator, so they descend to the quotient; we evaluate them on horizontal representatives.

\begin{definition}[Ambient potential on representatives]\label{def:rtc-ambient}
\lean{SymplecticPacking.Return.PhiE,SymplecticPacking.Return.lamE,SymplecticPacking.Return.OmegaE}
On $(\C^3\setminus0)\times\C$ set
\[
 q(Z,W)=\frac{|W|^2}{|Z|^6},\quad
 \Phi_E^{(\delta)}=\log|Z|^2+\delta\log(1+q),\quad
 \Omega_E^{(\delta)}=dd^c\Phi_E^{(\delta)},\quad \Omega_E=\Omega_E^{(d_n)}.
\]
On unit representatives put $p_W=d_n|W|^2/(1+|W|^2)$ and
$X_T=\{(Z,W):|Z|=1,\ \mathcal F(Z)=TW\}$.
Use $E_3$ as shorthand for these representative data with gauge $(Z,W)\sim(cZ,c^3W)$.
For $W\ne0$, write $\alpha=d\arg W/(2\pi)-3\eta_H$.
\end{definition}

A unit phase $e^{i\chi}$ shifts $\eta_H$ by $d\chi/(2\pi)$ and the fibre angular form by $3d\chi/(2\pi)$. Hence $\alpha$ is gauge invariant and vanishes on the gauge generator.

\begin{lemma}[$d^c$ of the norm]\label{lem:rtc-dc}
\lean{SymplecticPacking.Return.dc_log_normW, SymplecticPacking.Return.dc_log_one_add}
\uses{def:rtc-ambient}
On $\{W\ne0\}$, $d^c\log\|W\|^2=\alpha$. On $E_3$,
\[
 d^c\log\big(1+\|W\|^2\big)=\frac{\|W\|^2}{1+\|W\|^2}\,\alpha ,
\]
where the right side extends smoothly across $\{W=0\}$.
\end{lemma}

\begin{proof}
Choose a holomorphic lift $\hat Z$ and a fibre coordinate $\upsilon$ with $[Z,W]=[\hat Z,\upsilon]$. Then $\|W\|^2=|\upsilon|^2\|\hat Z\|^{-6}$ and $\arg W=\arg\upsilon$. Moreover $\eta_H=d^c\log\|\hat Z\|^2$, since $\operatorname{Re}(\hat Z^\dagger d\hat Z)=\frac12d\|\hat Z\|^2$. Hence $d^c\log\|W\|^2=d\arg\upsilon/2\pi-3\,d^c\log\|\hat Z\|^2=\alpha$.

The second identity follows from $d^c\phi(\mu)=\phi'(\mu)\mu\,d^c\log\mu$ with $\mu=\|W\|^2$. Smoothness follows from $\|W\|^2\alpha=\|\hat Z\|^{-6}[(\operatorname{Re}\upsilon\,d\operatorname{Im}\upsilon-\operatorname{Im}\upsilon\,d\operatorname{Re}\upsilon)/2\pi-3|\upsilon|^2d^c\log\|\hat Z\|^2]$.
\end{proof}

\begin{lemma}[The horizontal Hermitian blocks]\label{lem:rtc-kahler}
\lean{SymplecticPacking.Return.OmegaE_closed,SymplecticPacking.Return.OmegaE_herm,SymplecticPacking.Return.OmegaE_pos}
At $|Z|=1$, for $(\xi,\varsigma)\in H_Z\times\C$,
\begin{equation}\label{eq:rtc-OmegaE}
 \Omega_E^{(\delta)}((\xi_1,\varsigma_1),(\xi_2,\varsigma_2))
 =\frac1\pi\operatorname{Im}(A_\delta\xi_1^*\xi_2+B_\delta\bar\varsigma_1\varsigma_2),\quad
 A_\delta=1-\frac{3\delta|W|^2}{1+|W|^2},\quad B_\delta=\frac{\delta}{(1+|W|^2)^2}.
\end{equation}
For $\delta=d_n$, $A>1-3d_n>0$ and $B>0$.
The form is closed, and its restriction to every complex horizontal constraint kernel
$\{(\xi,\varsigma):d\mathcal F_Z\xi=T\varsigma\}$ is nondegenerate.
\end{lemma}

\begin{proof}
Since $q$ is $(\lambda Z,\lambda^3W)$-invariant for $\lambda\in\C^*$, it suffices to work on $H_Z\times\C$ at $|Z|=1$. There $\partial q=\bar W\,dW$, and
$\partial\bar\partial\log(1+q)=\frac{dW\wedge d\bar W}{(1+|W|^2)^2}-\frac{3|W|^2}{1+|W|^2}\sum_jdZ_j\wedge d\bar Z_j$, with no mixed terms. Applying $dd^c=\frac i{2\pi}\partial\bar\partial$ and adding $\omega_{\rm FS}$ gives \eqref{eq:rtc-OmegaE}. For $\delta=d_n$ the Hermitian form $A\xi_1^*\xi_2+B\bar\varsigma_1\varsigma_2$ is positive definite, so its imaginary part is nondegenerate on every complex subspace. Closedness is $d\,dd^c=0$.
\end{proof}

\subsection{The central normal injection}\label{sec:rtc-injection}

\begin{construction}[Weight-three injection]\label{cons:rtc-injection}
\lean{SymplecticPacking.Return.injection}
\uses{def:rtc-ambient, def:cub-model}
For a unit pair $(Z,U)$ with $p<d_n$ put
\[
 \iota[Z,U]=[Z,W],\qquad W=\frac{\sqrt\pi\,U}{\sqrt{d_n-p}},\qquad p=\pi|U|^2 .
\]
\end{construction}

\begin{lemma}[The fibre injection]\label{lem:rtc-injection}
\lean{SymplecticPacking.Return.injection_bijective,SymplecticPacking.Return.injection_normW,SymplecticPacking.Return.injection_pW,SymplecticPacking.Return.injection_normW_le}
At each unit $Z$, the radial maps
\[
 U\mapsto\frac{\sqrt\pi\,U}{\sqrt{d_n-\pi|U|^2}},\qquad
 W\mapsto W\sqrt{\frac{d_n}{\pi(1+|W|^2)}}
\]
are smooth inverses from $\{\pi|U|^2<d_n\}$ to $\C$ and back.
They commute with the weight-three gauge action and preserve $Z$.
Moreover $|W|^2=p/(d_n-p)$, $p_W=p$, and on $\mathcal N_h$ one has $|W|^2\le h/(d_n-h)$.
\end{lemma}

\begin{proof}
Substitute each radial expression into the other and use positivity of its denominator.
The norm and moment identities follow by squaring the positive radial factor.
The last bound follows from the increasing function $p\mapsto p/(d_n-p)$.
\end{proof}

\begin{proposition}[The central pullback on the cover]\label{prop:rtc-central}
\lean{SymplecticPacking.Return.injection_pullback}
On $\{p<d_n\}$ in the torus cover,
$(\iota\circ\mathsf N)^*\Omega_E=\Omega_0$.
\end{proposition}

\begin{proof}
By Lemma~\ref{lem:rtc-dc}, $d^c\Phi_E=\eta_H+p_W\alpha$. Along $\iota$ we have $p_W=p$ and $\arg W=\arg U$ (Lemma~\ref{lem:rtc-injection}). Along $\mathsf N$ we have $Z^*\eta_H=\eta$ and $U=w$. Hence $(\iota\circ\mathsf N)^*d^c\Phi_E=(1-3p)\eta+\frac12\operatorname{Im}(\bar w\,dw)$, which is smooth across $w=0$ and has exterior derivative $\Omega_0$ by Lemma~\ref{lem:fer-Omega0}.
\end{proof}

\subsection{Transport to $T=1$}\label{sec:rtc-transport}

The transport field is the lift of $\partial_T$ that preserves the constraint and is $\Omega_E$-orthogonal to $TX_T=\ker d(\mathcal F(Z)-TW)$.

\begin{construction}[Transport field]\label{cons:rtc-transport}
\lean{SymplecticPacking.Return.nvec, SymplecticPacking.Return.Gfun, SymplecticPacking.Return.Nden,
 SymplecticPacking.Return.transportField}
\uses{lem:rtc-kahler}
For $(Z,W,T)\in S^5\times\C\times\R$ put $f=\mathcal F(Z)$, $\bar Z^2=(\bar Z_0^2,\bar Z_1^2,\bar Z_2^2)$,
\[
 n=3\bar Z^2-3\bar fZ,\qquad G=\|n\|^2,\qquad N_{\rm den}=\frac GA+\frac{T^2}B ,
\]
with $A,B$ as in Lemma~\ref{lem:rtc-kahler}. On the open set $\mathcal O=\{N_{\rm den}>0\}$ define
\begin{equation}\label{eq:rtc-transport}
 \dot Z=\frac{W\,n}{A\,N_{\rm den}},\qquad \dot W=-\frac{T\,W}{B\,N_{\rm den}},\qquad \dot T=1 .
\end{equation}
\end{construction}

Since $A,B>0$, $N_{\rm den}\ge0$ with equality only where $T=0$ and $n=0$; so $\mathcal O$ is the complement of $\{T=0,\ n=0\}$, on which the field is smooth.

\begin{lemma}[Properties of the transport field]\label{lem:rtc-transport-field}
\lean{SymplecticPacking.Return.n_horizontal, SymplecticPacking.Return.n_dF, SymplecticPacking.Return.G_eq,
 SymplecticPacking.Return.transport_gaugeEquiv, SymplecticPacking.Return.transport_horizontal,
 SymplecticPacking.Return.transport_constraint, SymplecticPacking.Return.transport_normW,
 SymplecticPacking.Return.transport_orthogonal, SymplecticPacking.Return.transport_isotropic}
\uses{cons:rtc-transport, lem:rtc-kahler}
\begin{enumerate}[label=(\roman*)]
\item $Z^\dagger n=0$, $n^\dagger\xi=d\mathcal F_Z(\xi)$ for $\xi\in H_Z$, and $G=9\sum_j|Z_j|^4-9|f|^2$.
\item The field is equivariant under $(Z,W)\mapsto(e^{i\chi}Z,e^{3i\chi}W)$ and satisfies $Z^\dagger\dot Z=0$; hence it is
tangent to $S^5\times\C\times\R$. Its horizontal component is $V=(\dot Z,\dot W)\in H_Z\times\C$.
\item Along the field, $\frac{d}{dT}\mathcal F(Z)=\frac{WG}{AN_{\rm den}}=\frac{d}{dT}(TW)$, so $\mathcal F(Z)-TW$ is constant;
$\|Z\|$ is constant; and $\frac{d}{dT}|W|^2=-\frac{2T|W|^2}{BN_{\rm den}}\le0$ for $T\ge0$.
\item At a point of $X_T\times\{T\}$ in $\mathcal O$, $V$ is $\Omega_E$-orthogonal to $\ker(d\mathcal F_Z,-T)\subset H_Z\times\C$.
\end{enumerate}
\end{lemma}

\begin{proof}
(i) Since $Z^\dagger(3\bar Z^2)=3\bar f$, $n$ is the orthogonal projection of $3\bar Z^2$ to $H_Z$. This gives $Z^\dagger n=0$ and $G=\|3\bar Z^2\|^2-|3\bar f|^2$. For $\xi\in H_Z$, $d\mathcal F_Z(\xi)=(3\bar Z^2)^\dagger\xi=n^\dagger\xi$.

(ii) Under the gauge action $f\mapsto e^{3i\chi}f$ and $n\mapsto e^{-2i\chi}n$, while $G,A,B,N_{\rm den}$ are invariant. Tangency follows from (i). Finally, $H_Z\times\C$ is a complement to the kernel $\langle(iZ,3iW)\rangle$ of $T(S^5\times\C)\to TE_3$.

(iii) By (i), $\frac{d}{dT}\mathcal F(Z)=n^\dagger\dot Z$ and $\frac{d}{dT}(TW)=W+T\dot W$. The remaining claims follow from $\frac{d}{dT}\|Z\|^2=2\operatorname{Re}(Z^\dagger\dot Z)$ and $\frac{d}{dT}|W|^2=2\operatorname{Re}(\bar W\dot W)$.

(iv) By (i), $\ker(d\mathcal F_Z,-T)\cap(H_Z\times\C)$ is the kernel of $\ell(\xi,\varsigma)=n^\dagger\xi-T\varsigma$. Let $\mathfrak h=A\,\xi_1^\dagger\xi_2+B\,\bar\varsigma_1\varsigma_2$ and $\nu_\ell=(n/A,-T/B)$. Then $\mathfrak h(\nu_\ell,\kappa)=\ell(\kappa)$. Since $V=(W/N_{\rm den})\nu_\ell$ and $\mathfrak h$ is conjugate linear in its first slot, $\mathfrak h(V,\kappa)=0$ on $\ker\ell$, so $\Omega_E(V,\kappa)=\frac1\pi\operatorname{Im}\mathfrak h(V,\kappa)=0$. Also $\ell(V)=W$. Hence the field is the unique $\Omega_E$-orthogonal lift of $\partial_T$ preserving the constraint.
\end{proof}

\begin{lemma}[Denominator]\label{lem:rtc-denominator}
\lean{SymplecticPacking.Return.Nden_pos, SymplecticPacking.Return.Nden_lower}
\uses{cons:rtc-transport, lem:rtc-transport-field}
Let $\mathcal C=\{(Z,W,T):\|Z\|=1,\ \mathcal F(Z)=TW,\ 0\le T\le1\}$. On $\mathcal C$, $N_{\rm den}>0$: for $T=0$ one has $G\ge3$
and $N_{\rm den}\ge3$, and for $T>0$, $N_{\rm den}\ge T^2/B>0$. If moreover $|W|\le W_{\max}$ and
$T_0=\min\big(1,1/(\sqrt6\max(1,W_{\max}))\big)$, then $N_{\rm den}\ge\min\big(\frac32,\,T_0^2/d_n\big)$.
\end{lemma}

\begin{proof}
By Cauchy--Schwarz $\sum_j|Z_j|^4\ge\frac13$; since $0<A\le1$, $N_{\rm den}\ge G$; and $|f|=T|W|$ on $\mathcal C$. If $T=0$, $G=9\sum_j|Z_j|^4\ge3$. If $T\le T_0$, $9|f|^2\le9T_0^2W_{\max}^2\le\frac32$, so $G\ge\frac32$. If $T\ge T_0$, $N_{\rm den}\ge T^2/B\ge T_0^2/d_n$, since $B\le d_n$.
\end{proof}

\begin{proposition}[Protected transport]\label{prop:rtc-transport}
\lean{SymplecticPacking.Return.transport_mem,SymplecticPacking.Return.transport_hasDerivAt,SymplecticPacking.Return.transport_smooth,SymplecticPacking.Return.transport_injective,SymplecticPacking.Return.transport_pullback}
Fix $w_*=h/(d_n-h)$ and let $\tau_T$ be the flow, from $T=0$, of the cut field defined in the proof. It is a smooth
injective map, commuting with the gauge action. For every unit pair $(Z,W)\in X_0$ with $|W|^2\le w_*$ the trajectory
solves \eqref{eq:rtc-transport} on $[0,1]$, stays in $X_T\cap\{|W|^2\le w_*\}$, and $\tau_1$ maps
$X_0\cap\{|W|^2\le w_*\}$ into $X_1$ with $\tau_1^*(\Omega_E|_{X_1})=\Omega_E|_{X_0}$ there.
\end{proposition}

\begin{proof}
\emph{Cut field.} Let $T_0=T_0(\sqrt{w_*})$ and $c_*=\min(\frac32,T_0^2/d_n)$. For $a<b$ write $S_{a,b}(x)=S((x-a)/(b-a))$ and $P_{a,b,c,d}(x)=S_{a,b}(x)(1-S_{c,d}(x))$, and put
\[
 \chi=S_{c_*/2,3c_*/4}(N_{\rm den})\,P_{1/2,3/4,5/4,3/2}(|Z|^2)\,S_{0,1}(w_*+2-|W|^2)\,P_{-1,-1/2,3/2,2}(T).
\]
Off the unit sphere, $n$ is replaced by $3\bar Z^2-3\bar fZ/|Z|^2$ in computing $N_{\rm den}$. The field $\chi\cdot(\dot Z,\dot W,1)$ is smooth, since its first factor vanishes where $N_{\rm den}\le c_*/2$. It is supported in $\{\frac12\le|Z|^2\le\frac32,\ |W|^2\le w_*+2,\ |T|\le2\}$, so its flow is complete, with diffeomorphic time-$T$ maps. The flow is gauge equivariant because $\chi$ is gauge invariant and the field is equivariant (Lemma~\ref{lem:rtc-transport-field}(ii)).

\emph{Trajectories.} On $\mathcal C_*=\{|Z|=1,\ \mathcal F(Z)=TW,\ |W|^2\le w_*,\ 0\le T\le1\}$ we have $\chi=1$, by Lemma~\ref{lem:rtc-denominator} with $W_{\max}=\sqrt{w_*}$. By Lemma~\ref{lem:rtc-transport-field}(iii), uncut trajectories from $X_0\cap\{|W|^2\le w_*\}$ stay in $\mathcal C_*$, so by uniqueness they coincide with cut trajectories on $[0,1]$.

\emph{Pullback.} Let $F(\sigma,T)=\tau_T(x(\sigma))$ for a smooth family $x(\sigma)$ of initial points. Since $\mathcal F(Z)-TW\equiv0$ along $F$, each tangent vector $(\kappa,t)$ satisfies $\ell(\kappa)=tW$, where $\ell=(d\mathcal F_Z,-T)$. Replacing $\kappa$ by its horizontal representative $\kappa_H$ modulo $G_{\rm gauge}=(iZ,3iW)$ changes neither the constraint, since $\ell(G_{\rm gauge})=3i(\mathcal F(Z)-TW)=0$, nor the pairing with $\Omega_E$. Now $\partial_TF=(V,1)$ with $\ell(V)=W$, so $\kappa_H-tV\in\ker\ell$. Lemma~\ref{lem:rtc-transport-field}(iv) then gives $\Omega_E(V,\kappa)=0$, i.e.\ $\iota_{\partial_T}F^*\Omega_E=0$. Since $F^*\Omega_E$ is closed, Cartan's formula gives $\partial_T(F_T^*\Omega_E)=0$, hence $F_1^*\Omega_E=F_0^*\Omega_E$.
\end{proof}

\subsection{Return to the Fubini--Study form}\label{sec:rtc-return}

On $X_1$ we have $W=\mathcal F(Z)$, so $\mathrm{pr}$ has inverse $\mathrm{gr}[Z]=[Z,\mathcal F(Z)]$. Let $u=|\mathcal F(Z)|^2/\|Z\|^6$, which is smooth on unit representatives. For $s\in\R$ put $\omega^{\rm gr}_s=\omega_{\rm FS}+s\,d_n\,dd^c\log(1+u)=\mathrm{gr}^*\Omega_E^{(sd_n)}$, where the equality uses $\mathrm{gr}^*\|W\|^2=u$. Then $\omega^{\rm gr}_0=\omega_{\rm FS}$ and $\omega^{\rm gr}_1=\mathrm{gr}^*\Omega_E=(\mathrm{pr}^{-1})^*(\Omega_E|_{X_1})$. The return is the Moser flow of this path, which requires $\iota_{Y_s}\omega^{\rm gr}_s=-d_n\,d^c\log(1+u)$; this is solved in closed form because $n$ is an eigenvector of $M_s$.

\begin{construction}[Return field]\label{cons:rtc-return}
\lean{SymplecticPacking.Return.omegaS, SymplecticPacking.Return.returnField}
\uses{lem:rtc-transport-field}
In unit representatives, with $f$, $n$, $G$ as in Construction~\ref{cons:rtc-transport}, put
\[
 A_s=1-3s\,d_n\frac{u}{1+u},\qquad B_s=\frac{s\,d_n}{(1+u)^2},\qquad M_s=A_sI+B_s\,nn^\dagger\ \text{ on }H_Z,
\]
\begin{equation}\label{eq:rtc-return-field}
 Y_s=-\frac{d_n}{2}\,\frac{f\,n}{(1+u)(A_s+B_sG)}\in H_Z .
\end{equation}
\end{construction}

\begin{lemma}[Graph forms]\label{lem:rtc-graph}
\lean{SymplecticPacking.Return.omegaS_graph, SymplecticPacking.Return.OmegaE_herm,
 SymplecticPacking.Return.As_gt_half, SymplecticPacking.Return.Bs_nonneg, SymplecticPacking.Return.returnDen_gt_half, SymplecticPacking.Return.Ms_nvec}
\uses{cons:rtc-return, lem:rtc-kahler}
For $s\in[0,1]$ and $\xi_1,\xi_2\in H_Z$, $\omega^{\rm gr}_s(\xi_1,\xi_2)=\frac1\pi\operatorname{Im}(\xi_1^\dagger M_s\xi_2)$. Moreover
$u\le1$, $A_s\ge1-\frac32d_n>\frac12$ and $B_s\ge0$; hence $M_s\ge\frac12I$ and every $\omega^{\rm gr}_s$ is positive on $H_Z$.
\end{lemma}

\begin{proof}
By Lemma~\ref{lem:rtc-transport-field}(i), $d\,\mathrm{gr}(\xi)$ is represented by $(\xi,n^\dagger\xi)$ in the gauge of Lemma~\ref{lem:rtc-kahler}. Since $|W|^2=u$ at $W=\mathcal F(Z)$, \eqref{eq:rtc-OmegaE} gives the Hermitian form $A_s\xi_1^\dagger\xi_2+B_s\overline{n^\dagger\xi_1}\,n^\dagger\xi_2=\xi_1^\dagger M_s\xi_2$. For $\|Z\|=1$, $|f|\le\max_j|Z_j|\sum_j|Z_j|^2\le1$. Hence $u\le1$, so $u/(1+u)\le\frac12$ and $A_s\ge1-\frac32d_n>\frac12$, because $d_n<\frac13$.
\end{proof}

\begin{lemma}[$d^c$ along the graph]\label{lem:rtc-dc-horizontal}
\lean{SymplecticPacking.Return.dc_log_one_add_u}
\uses{cons:rtc-return, lem:rtc-transport-field}
For $\xi\in H_Z$ (unit representative), $d^c\log(1+u)(\xi)=\dfrac{\operatorname{Im}\big(\bar f\,n^\dagger\xi\big)}{2\pi(1+u)}$.
\end{lemma}

\begin{proof}
Where $f\ne0$, as in Lemma~\ref{lem:rtc-dc}, $d^c\log(1+u)=\tfrac{u}{1+u}\big(\tfrac{d\arg\mathcal F}{2\pi}-3\eta_H\big)$. On $H_Z$ we have $\eta_H(\xi)=0$ and $d\arg\mathcal F(\xi)=\operatorname{Im}(n^\dagger\xi/f)$, and $u/f=\bar f$. Both sides are smooth, so the identity extends to $\{f=0\}$.
\end{proof}

\begin{lemma}[The return field]\label{lem:rtc-return-field}
\lean{SymplecticPacking.Return.returnField_spec, SymplecticPacking.Return.returnField_gauge}
\uses{lem:rtc-graph, lem:rtc-dc-horizontal}
$Y_s$ is smooth in $(s,Z)$ for $s$ in a neighbourhood of $[0,1]$, gauge equivariant and tangent to $S^5$, hence descends to a
time-dependent smooth vector field on unit representatives, and
\[
 \iota_{Y_s}\omega^{\rm gr}_s=-d_n\,d^c\log(1+u)\qquad(s\in[0,1]).
\]
\end{lemma}

\begin{proof}
Under the gauge action $f\mapsto e^{3i\chi}f$ and $n\mapsto e^{-2i\chi}n$, so $Y_s\mapsto e^{i\chi}Y_s$; also $Z^\dagger Y_s=0$. The denominator is positive for $-\frac1{9d_n}<s<\frac2{3d_n}$, an interval containing $[0,1]$: for $s\ge0$, $A_s+B_sG\ge A_s\ge1-\frac32s\,d_n$; for $s<0$, $A_s\ge1$ and $B_sG\ge-9|s|d_n$, because $0\le G\le9$. Since $M_sn=(A_s+B_sG)n$, we get $M_sY_s=-\frac{d_n}2\frac{fn}{1+u}$. Lemmas~\ref{lem:rtc-graph} and~\ref{lem:rtc-dc-horizontal} then give, for $\xi\in H_Z$, $\omega^{\rm gr}_s(Y_s,\xi)=\frac1\pi\operatorname{Im}((M_sY_s)^\dagger\xi)=-d_n\,d^c\log(1+u)(\xi)$.
\end{proof}

\begin{proposition}[Unit-representative return]\label{prop:rtc-return}
\lean[needed]{SymplecticPacking.Return.returnFlow_pullback,SymplecticPacking.Return.returnFlowInv_pullback,
SymplecticPacking.Return.returnFlow_smooth,SymplecticPacking.Return.returnFlowInv_smooth,
SymplecticPacking.Return.returnFlowInv_returnFlow,SymplecticPacking.Return.returnFlow_returnFlowInv,
SymplecticPacking.Return.nsq_returnFlow,SymplecticPacking.Return.returnFlow_gaugeEquiv}
Let $\psi_s$ be the flow from $s=0$ of the cut field of the proof. It is defined for all $s$, preserves $|Z|=1$, commutes
with the gauge action, and its time-$s$ maps are diffeomorphisms. On the unit sphere it is the flow of $Y_s$ for
$s\in[0,1]$, and $\psi_s^*\omega_s^{\rm gr}=\omega_{\rm FS}$. Consequently $(\psi_1^{-1})^*\omega_{\rm FS}=\omega^{\rm gr}_1$, and
$(\psi_1^{-1}\circ\mathrm{pr})^*\omega_{\rm FS}=\Omega_E|_{X_1}$.
\end{proposition}

\begin{proof}
Put $m=1-\frac32d_n>0$ and $D_s=A_s+B_sG$, with $n$ extended off the sphere as in Proposition~\ref{prop:rtc-transport}. The cut field is
\[
 S_{m/2,3m/4}(D_s)\,P_{1/2,3/4,5/4,3/2}(|Z|^2)\,P_{-1,-1/2,3/2,2}(s)\,Y_s .
\]
It is smooth, since the first factor vanishes where $D_s\le m/2$, and compactly supported, so its flow is complete with diffeomorphic time maps. The flow is gauge equivariant by Lemma~\ref{lem:rtc-return-field}, and it preserves the sphere because $2\operatorname{Re}(Z^\dagger Y_s)=0$. For $|Z|=1$ and $s\in[0,1]$, Lemma~\ref{lem:rtc-graph} gives $D_s\ge m$, so the cut field equals $Y_s$ there. Since $\partial_s\omega^{\rm gr}_s=d(d_n\,d^c\log(1+u))$, Cartan's formula and Lemma~\ref{lem:rtc-return-field} give $\frac d{ds}\psi_s^*\omega^{\rm gr}_s=0$. Finally, $\mathrm{gr}\circ\mathrm{pr}=\mathrm{id}$ on $X_1$ gives $\mathrm{pr}^*\omega^{\rm gr}_1=\Omega_E|_{X_1}$.
\end{proof}

\subsection{Equivariant injectivity and the ball chart}\label{sec:rtc-avoid}

From now on, $\mathcal T_{\rm u}$ uses the protected cutoff $w_*=h/(d_n-h)$ and is considered on $\mathcal N_h$. Let $\mathsf r Z=(\zeta Z_0,Z_1,Z_2)$ and $\mathsf r(Z,W)=(\mathsf r Z,W)$; $\mathsf r$ preserves the norm, the Fermat equation and the source moment, and commutes with the unit gauge action.

\begin{lemma}[Fixed representatives modulo gauge]\label{lem:rtc-fixed}
\lean{SymplecticPacking.Return.fixed_CP2,SymplecticPacking.Return.fixed_N}
For a nonzero vector, $\mathsf r Z$ is a scalar multiple of $Z$ exactly when $Z_0=0$ or $Z_1=Z_2=0$.
For a central unit pair $(Z,U)$, the equation $\mathsf r(Z,U)=(cZ,c^3U)$ is possible exactly when $Z_0=0$.
\end{lemma}

\begin{proof}
A nonzero last coordinate forces $c=1$ and then $Z_0=0$. Otherwise both last coordinates vanish, which $\mathcal F(Z)=0$ excludes at a central unit vector. If $Z_0=0$, the pair is fixed with $c=1$.
\end{proof}

\begin{lemma}[Equivariance]\label{lem:rtc-equivariant}
\lean{SymplecticPacking.Return.injection_equivariant, SymplecticPacking.Return.transport_equivariant,
 SymplecticPacking.Return.returnFlow_equivariant, SymplecticPacking.Return.compositeT_equivariant}
\uses{cons:rtc-injection, prop:rtc-transport, prop:rtc-return}
$\mathsf r$ preserves $\mathcal F$, $\omega_{\rm FS}$, $\|W\|$, $\Omega_E$, $u$, each $X_T$ and $\{(Z,U):|Z|=1,\mathcal F(Z)=0,p<d_n\}$, and it maps the fields
\eqref{eq:rtc-transport} and \eqref{eq:rtc-return-field} to themselves. Consequently $\iota$, $\tau_1$, $\mathrm{pr}$, $\psi_1^{-1}$ and
$\mathcal T_{\rm u}$ commute with $\mathsf r$.
\end{lemma}

\begin{proof}
$\mathsf r$ is unitary and $\mathcal F\circ\mathsf r=\mathcal F$, so it preserves all the listed quantities and also $p,A,B,A_s,B_s$. Since $\bar\zeta^2=\zeta$, we have $n(\mathsf rZ)=\mathsf r\,n(Z)$. Hence $G$ and $N_{\rm den}$ are invariant, both fields are $\mathsf r$-equivariant, and by uniqueness so are their flows. $\iota$ commutes with $\mathsf r$ because $p$ is invariant, and $\mathrm{pr}$ does because the lift of $\mathsf r$ to $E_3$ covers $\mathsf r$.
\end{proof}

\begin{lemma}[Injectivity modulo gauge]\label{lem:rtc-injective}
\lean{SymplecticPacking.Return.compositeT_inj_mod_gauge}
For $x,x'\in\mathcal N_h$ and $|c|=1$, the equality
$\mathcal T_{\rm u}(x)=c\mathcal T_{\rm u}(x')$ implies
$x=(cZ',c^3U')$.
\end{lemma}

\begin{proof}
Each factor of $\mathcal T_{\rm u}=\psi_1^{-1}\circ\mathrm{pr}\circ\tau_1\circ\iota$ is injective on representatives: $\iota$ by the explicit inverse of Lemma~\ref{lem:rtc-injection}, $\tau_1$ and $\psi_1^{-1}$ as time maps of complete flows, and $\mathrm{pr}$ on $X_1$, where $W=\mathcal F(Z)$. Each factor also commutes with $(Z,W)\mapsto(cZ,c^3W)$. Hence $c\,\mathcal T_{\rm u}(x')=\mathcal T_{\rm u}(cZ',c^3U')$, and injectivity concludes.
\end{proof}

\begin{proposition}[Avoidance and compact margin]\label{prop:rtc-avoid}
\lean{SymplecticPacking.Return.avoid_L,SymplecticPacking.Return.avoid_L_compact}
If $x=(Z,U)\in\mathcal N_h$ and $Z_0\ne0$, then
$(\mathcal T_{\rm u}(x))_0\ne0$ and its last two coordinates do not vanish together.
On every nonempty compact family of these inputs,
$m_{\mathcal K}=\min_{x\in\mathcal K}|(\mathcal T_{\rm u}(x))_0|^2>0$. For $\mathcal K=\varnothing$ set $m_{\mathcal K}=1$.
\end{proposition}

\begin{proof}
Let $Y=\mathcal T_{\rm u}(x)$. If $Y_0=0$, then $\mathsf rY=Y$, and Lemmas~\ref{lem:rtc-equivariant} and~\ref{lem:rtc-injective} (with $c=1$) give $\mathsf rx=x$. If $Y_1=Y_2=0$, then $\mathsf rY=\zeta Y$, and the same lemmas give $\mathsf rx=(\zeta Z,\zeta^3U)$. In both cases Lemma~\ref{lem:rtc-fixed} gives $Z_0=0$, a contradiction. The compact margin holds because $x\mapsto|(\mathcal T_{\rm u}(x))_0|^2$ is continuous and positive on $\mathcal K$.
\end{proof}

\begin{lemma}[The ball chart on representatives]\label{lem:rtc-ball-chart}
\lean[needed]{SymplecticPacking.Return.ballChart, SymplecticPacking.Return.ballChart_gauge,
 SymplecticPacking.Return.ballChart_symplectic_rep, SymplecticPacking.Return.ballChart_affine,
 SymplecticPacking.Return.ballChart_norm, SymplecticPacking.Return.ballAffine_eq, SymplecticPacking.Return.ballChart_norm_lt}
On $S^5_*=\{Z\in\C^3:|Z|=1,\ Z_0\ne0\}$ put $b_{\rm ball}(Z)=\frac{\bar Z_0}{|Z_0|}\frac{(Z_1,Z_2)}{\sqrt\pi}$. This map is gauge
invariant, $\pi|b_{\rm ball}(Z)|^2=1-|Z_0|^2<1$, and in the affine coordinates $z=(Z_1/Z_0,Z_2/Z_0)$ it is the chart
$\mathrm{ch}(z)=z/\sqrt{\pi(1+|z|^2)}$ of Lemma~\ref{lem:pre-fs}(b). Moreover
$b_{\rm ball}^*\omega_0=\frac1\pi\sum_{j=0}^2dx_j\wedge dy_j\big|_{S^5_*}=\omega_{\rm FS}|_{S^5_*}$, where $Z_j=x_j+iy_j$. Hence
$(b_{\rm ball}\circ\varphi)^*\omega_0=\varphi^*\omega_{\rm FS}$ for every smooth map $\varphi$ of a manifold into $S^5_*$.
\end{lemma}

\begin{proof}
The unit phases cancel, and substituting $Z=e^{i\chi}(1,z)/\sqrt{1+|z|^2}$ gives the affine formula and the norm identity.
Write $e=\bar Z_0/|Z_0|=e^{i\phi}$, $\phi=-\arg Z_0$, and $y_j=eZ_j/\sqrt\pi$ ($j=1,2$). Then
$\bar y_jdy_j=\frac1\pi\big(|Z_j|^2\,\bar e\,de+\bar Z_jdZ_j\big)$ with $\bar e\,de=i\,d\phi$, so the Liouville form
$\lambda_0=\frac12\sum_{j=1}^2\operatorname{Im}(\bar y_jdy_j)$ pulls back to
$\frac1{2\pi}\big((1-|Z_0|^2)d\phi+\sum_{j=1}^2\operatorname{Im}(\bar Z_jdZ_j)\big)$ on $S^5_*$. Since
$\operatorname{Im}(\bar Z_0dZ_0)=|Z_0|^2d\arg Z_0=-|Z_0|^2d\phi$, this equals $\frac1{2\pi}\big(d\phi+\operatorname{Im}(Z^\dagger dZ)\big)$, whose
differential is $\frac1{2\pi}\sum_{j=0}^2d(x_jdy_j-y_jdx_j)=\frac1\pi\sum_jdx_j\wedge dy_j$. As $d\lambda_0=\omega_0$, this proves the first
identity; the second holds because $\frac1{2\pi}\operatorname{Im}(Z^\dagger dZ)=\eta_H$ on the sphere and $d\eta_H=\omega_{\rm FS}$.
\end{proof}

For $y\in\Bo(1)$, $(\sqrt{1-\pi|y|^2},\sqrt\pi\,y_1,\sqrt\pi\,y_2)$ is the unique gauge representative with positive first coordinate of every unit vector with nonzero first coordinate and ball-chart image $y$.

\subsection{The transfer theorem}\label{sec:rtc-transfer}

\begin{theorem}[Transfer on the cover]\label{thm:rtc-transfer}\label{thm:cub-transfer}
\lean[needed]{SymplecticPacking.Return.compositeT_inj_mod_gauge,SymplecticPacking.Return.compositeT_pullback,SymplecticPacking.Return.compositeT_equivariant,SymplecticPacking.Return.avoid_L_compact,
SymplecticPacking.Return.injection_pullback, SymplecticPacking.Return.ballChart_symplectic_rep,
SymplecticPacking.Embed.Cubic.fcap_eq_imp, SymplecticPacking.Embed.Cubic.fcap_eq_imp}
\uses{cons:cub-transfer, prop:rtc-central, prop:rtc-transport, prop:rtc-return, prop:rtc-avoid, lem:rtc-ball-chart}
Let $w_*=h/(d_n-h)$ and $\mathcal T_{\rm u}=\psi_1^{-1}\circ\mathrm{pr}\circ\tau_1\circ\iota$ on $\mathcal N_h$. The cover map
\[
 \mathcal T=b_{\rm ball}\circ\mathcal T_{\rm u}\circ\mathsf N
 \quad\text{on}\quad\{p<h,\ (u,v)\notin\mathcal Z\}
\]
is smooth, invariant under the deck maps, takes values in $\Bo(1)$ and satisfies $\mathcal T^*\omega_0=\Omega_0$. Two points
have the same image if and only if they are deck equivalent. For every compact $\mathcal K$ in the domain,
$\mathcal T(\mathcal K)\subset B(1-m_{\mathcal K})$ with $m_{\mathcal K}>0$ as in Proposition~\ref{prop:rtc-avoid}. In particular the
restriction of $\mathcal T$ to an open set on which it is injective is a symplectic embedding, and $\mathcal T$ carries
pairwise disjoint subsets of its domain, no two points of whose union are deck equivalent, to pairwise disjoint sets.
\end{theorem}

\begin{proof}
On $\{p<h\}$, $\mathsf N$ takes values in $\mathcal N_h$ with $p<d_n$, and $\tau_1\circ\iota$ lands in $X_1\cap\{|W|^2\le w_*\}$ (Lemma~\ref{lem:rtc-injection}, Proposition~\ref{prop:rtc-transport}). Propositions~\ref{prop:rtc-central}, \ref{prop:rtc-transport} and~\ref{prop:rtc-return} give $(\mathcal T_{\rm u}\circ\mathsf N)^*\omega_{\rm FS}=\Omega_0$. For $(u,v)\notin\mathcal Z$, $\mathcal T_{\rm u}(\mathsf N(x))\in S^5_*$ (Proposition~\ref{prop:rtc-avoid}), so Lemma~\ref{lem:rtc-ball-chart} gives $\mathcal T^*\omega_0=\Omega_0$ and $\pi|\mathcal T|^2=1-|(\mathcal T_{\rm u}\mathsf N)_0|^2<1$, whence the margin $m_{\mathcal K}$ on $\mathcal K$. Deck maps change $\mathsf N(x)$ to a gauge-equivalent pair (Construction~\ref{cons:fer-atlas}(b)), $\mathcal T_{\rm u}$ is gauge equivariant and $b_{\rm ball}$ gauge invariant, so $\mathcal T$ is deck invariant. Conversely, if $\mathcal T(x)=\mathcal T(x')$, the ball chart determines a unit vector with nonzero first coordinate up to gauge, so Lemma~\ref{lem:rtc-injective} gives $\mathsf N(x)\sim\mathsf N(x')$; hence $[Z(u,v)]=[Z(u',v')]$ in $\CP^2$, and Lemma~\ref{lem:fer-torus} gives $(u,v)-(u',v')\in\Z^2$. For the corresponding deck map $T$, the unit vectors $Z(u,v)$ agree, so the gauge factor is one and $x=Tx'$. Injective restrictions are injective symplectic immersions in equal dimension, hence open symplectic embeddings; the last assertion follows from the characterisation of equal images.
\end{proof}

All objects are explicit; $\psi_1^{-1}$ is obtained by integrating \eqref{eq:rtc-return-field} backwards from $s=1$, and the transport by \eqref{eq:rtc-transport}.

\subsection{The packing in dimension four}\label{sec:rtc-packing}

The transfer depends on the data only through $h$.

\begin{construction}[Data of the four-dimensional packing]\label{cons:rtc-data}
\uses{M-prop:cmp-linear, def:shd4-data, prop:cub-handle, prop:fer-handle, thm:rtc-transfer, prop:rtc-avoid}
Let $(k,\rho_i,a_i,h,d_n,g,K,q)$ be valid cubic data (Definition~\ref{def:shd4-data}) with $d_n=\frac12(h+\frac13)$. Let
$f_i=f_i^{\rm cap}$ be the maps of Remark~\ref{rem:shd-four}, built by the handle construction of \S\ref{sec:shd-four},
put $\bar\rho_i=(\rho_i+a_i)/2$ and $\hat\rho_i=\bar\rho_i+(a_i-\bar\rho_i)/4$, and let the open domain be
$\Bo(\hat\rho_i)$. Let $\rho_i^+<\rho_i^{\rm dom}$ be the recorded capacities of Remark~\ref{rem:rtc-recorded}. Put
$\mathcal K=\bigcup_i f_i(B(\rho_i^+))$ and $m=\frac12\min_{z\in\mathcal K}(1-\pi|z|^2)$.

\emph{The equal case.} Let $k\ge10$, $\rho_1=\dots=\rho_k=\rho$ and $\frac34<\rho\sqrt k<1$. Put
$t_0=\frac12(1+\rho\sqrt k)$, $a_i=a=t_0/\sqrt k$, $h=t_0/3$, $g=3(1-t_0)$, and let $(K,q)$ be the comparison datum of
\Mref{Proposition}{prop:cmp-linear}, with gap $\frac{31}{32}g$. These are valid cubic data
(Remark~\ref{rem:rtc-equal}). For $k=9$ the same formulas would give $a=t_0/3=h$, violating $a<h$; valid cubic data for
$k=9$ and for unequal capacities are constructed in the companion paper~\cite{SongHD} from a scalar comparison datum.
\end{construction}

\begin{theorem}[Packing in dimension four]\label{thm:cub-packing}\label{prop:rtc-embeddings}
\lean[needed]{SymplecticPacking.Branch2.CubicData.Valid.fcapGen_injective, SymplecticPacking.Embed.Cubic.fcapV_disjoint,
SymplecticPacking.Embed.Cubic.fcapV_mem_ball, SymplecticPacking.Embed.Cubic.fcap_injective, SymplecticPacking.Embed.Cubic.fcap_disjoint,
SymplecticPacking.Embed.Cubic.fcap_mem_ball, SymplecticPacking.Embed.Cubic.handle_phiInvSpec,
SymplecticPacking.Handle.Cubic.handle_deckEquivariant}
\uses{cons:rtc-data, M-prop:cmp-linear, def:shd4-data, prop:cub-handle, prop:fer-handle, thm:rtc-transfer, prop:rtc-avoid}
Let $(k,\rho_i,a_i,h,d_n,g,K,q)$ be valid cubic data with $d_n=\frac12(h+\frac13)$. Then there are explicit numbers
$\rho_i<\rho_i^+<\rho_i^{\rm dom}$ and smooth maps $f_i\colon\Bo(\rho_i^{\rm dom})\to\C^2$ with symplectic derivative, injective
on $B(\rho_i^+)$, such that the sets $f_i(B(\rho_i^+))$ are pairwise disjoint and contained in $B(1-\mu_{\rm ball})$ for some
$\mu_{\rm ball}\in(0,1)$. More precisely, for the data of Construction~\ref{cons:rtc-data}:
\begin{enumerate}[label=(\roman*)]
\item each $f_i$ is a smooth symplectic embedding of $\Bo(\hat\rho_i)$ into $(\Bo(1),\omega_0)$, and
$\rho_i<\rho_i^+<\rho_i^{\rm dom}\le\hat\rho_i$;
\item the sets $f_i(\Bo(\hat\rho_i))$, $1\le i\le k$, are pairwise disjoint;
\item $m>0$ and $f_i(B(\rho_i^+))\subset B(1-m)$ for every $i$;
\item the comparison isotopy $q$, the Moser map $\Phi$, the handle map $\mathcal D$ and the transfer $\mathcal T$ are the same for all $i$;
only the layer map $E_i$ depends on $i$.
\end{enumerate}
\end{theorem}

\begin{proof}
The data are valid cubic data by hypothesis (in the equal case by \Mref{Proposition}{prop:cmp-linear} and Remark~\ref{rem:rtc-equal}), and $h<d_n<\frac13$. Here $f_i=\mathcal T\circ f^N_i$ with $f^N_i$ of Proposition~\ref{prop:shd4-maps} (Remark~\ref{rem:shd-four}).

(i) By Proposition~\ref{prop:shd4-maps}(a), $f^N_i$ is smooth on $\Bo(\hat\rho_i)$ into $\{p<h\}\setminus\widetilde{\mathcal F}$, the domain of $\mathcal T$, with $(f^N_i)^*\Omega_0=\omega_0$; Theorem~\ref{thm:rtc-transfer} gives $f_i$ smooth into $\Bo(1)$ with $f_i^*\omega_0=\omega_0$. If $f_i(X)=f_i(X')$, Theorem~\ref{thm:rtc-transfer} gives $f^N_i(X)=Tf^N_i(X')$ for a deck map $T$, and Proposition~\ref{prop:shd4-maps}(c) gives $X=X'$; so $f_i$ is an open embedding. The capacity inequalities are in Remark~\ref{rem:rtc-recorded}.
(ii) If $f_i(X)=f_j(X')$, the same argument gives $i=j$, $X=X'$.
(iii) $\mathcal K$ is compact in $\Bo(1)$, so $1-\pi|z|^2\ge2m>m>0$ on $\mathcal K$.
(iv) $q$, $\mathcal D$, $\Phi$, $\chi_{\mathcal A}$ and $\mathcal T$ do not depend on $i$.

The first statement follows with $\mu_{\rm ball}=m<1$, restricting $f_i$ to $\Bo(\rho_i^{\rm dom})$.
\end{proof}

\begin{remark}[Recorded capacities]\label{rem:rtc-recorded}
Since the proof embeds the open balls $\Bo(\hat\rho_i)$ with pairwise disjoint images, Theorem~\ref{thm:cub-packing}
holds for every choice of $\rho_i<\rho_i^+<\rho_i^{\rm dom}\le\hat\rho_i$. The choice enters only the radii of the common
isotopy (\Mref{Definition}{def:ham-radii}) and hence the Hamiltonian. In the equal case the construction records
$(\rho_i^+,\rho_i^{\rm dom})=(\rho \varrho_3^2,\rho \varrho_4^2)$, where $\varrho_4=\sqrt{\bar\rho/\rho}$ and $\varrho_3=(1+\varrho_4)/2$; then
$1<\varrho_3<\varrho_4$ gives $\rho<\rho \varrho_3^2<\rho \varrho_4^2=\bar\rho<\hat\rho$. For target capacity $C$ and given capacity $c=C\rho$
the physical map $u\mapsto\sqrt C\,f_i(u/\sqrt C)$ has the capacities $C\rho_i^+$ and $C\rho_i^{\rm dom}$, that is,
$c\varrho_3^2$ and $c\varrho_4^2$. For the valid cubic data of the companion paper~\cite{SongHD} the construction records
$(\rho_i^+,\rho_i^{\rm dom})=(\bar\rho_i,\hat\rho_i)$.
\end{remark}

\begin{remark}[The equal-capacity branch]\label{rem:rtc-equal}\label{rem:cub-lean}
\lean{SymplecticPacking.Params.t0, SymplecticPacking.Params.a, SymplecticPacking.Params.h,
 SymplecticPacking.Params.dn_sub_h, SymplecticPacking.Params.one_sub_three_dn, SymplecticPacking.Params.a_sub_rho,
 SymplecticPacking.Params.density, SymplecticPacking.Branch2.CubicData.ofParams_valid,
 SymplecticPacking.Branch2.embeddingsEq, SymplecticPacking.Branch2.embeddingsEq_spec,
 SymplecticPacking.Branch2.cubicEmbeddings_spec}
\uses{cons:rtc-data, thm:cub-packing, rem:fer-lean}
In the equal case of Construction~\ref{cons:rtc-data} let $\eta_\star=1-\rho\sqrt k\in(0,\frac14)$ be the volume defect. Then
\[
 t_0=1-\frac{\eta_\star}2,\quad a=\frac{t_0}{\sqrt k},\quad h=\frac{t_0}3,\quad d_n=\frac13-\frac{\eta_\star}{12},\qquad
 a-\rho=\frac{\eta_\star}{2\sqrt k},\quad d_n-h=\frac{\eta_\star}{12},\quad 1-3d_n=\frac{\eta_\star}4,
\]
the comparison function is $K(p)=(3\sqrt k-9)p$, and the comparison gap of \Mref{Proposition}{prop:cmp-linear} is
$\frac{31}{32}g$ with $g=3(1-t_0)=\frac{3\eta_\star}2$. With the isotopy of Theorem~\ref{thm:pll-inner} and the handle
construction of \S\ref{sec:shd-four}, the maps of this branch are
\[
 f_i= b_{\rm ball}\circ\psi_1^{-1}\circ\operatorname{pr}\circ\tau_1\circ\iota
       \circ\mathsf N\circ\Phi^{-1}\circ\chi_{\mathcal A}\circ\mathcal D^{-1}\circ E_i\qquad(1\le i\le k).
\]
For the input of \Mref{Corollary}{thm:main-epoch}, $\rho=\sqrt{(1-\eps/2)/k}$ and
$\eta_\star=1-\sqrt{1-\eps/2}\in(\eps/4,\eps/2)$; the equal case applies exactly when $\rho\sqrt k>3/4$, i.e.\ $\eps<\frac78$.
For unit balls (capacity $\pi$) the target capacity is $C=\pi/\rho$, and the
physical rescaling of \Mref{Definition}{def:ham-normalise} makes the target the open Euclidean ball of radius
$R=\rho^{-1/2}$, so $k/R^4=k\rho^2=(1-\eta_\star)^2$. For the data of \Mref{Theorem}{thm:intro-4d}, $R^4=k/(1-\eps/4)$, so
$\rho\sqrt k=\sqrt{1-\eps/4}>\frac34$ and $\eta_\star=1-\sqrt{1-\eps/4}\in(\eps/8,\eps/4)$; the equal case applies
for every $\eps\in(0,1)$ when $k\ge10$. The strip constants are those of Remark~\ref{rem:shd-four} with
$\delta=(a-\rho)/2$.

The defect $1-\rho\sqrt k=\eta_\star$ splits as $(1-t_0)+\sqrt k\,(a-\rho)=\frac{\eta_\star}2+\frac{\eta_\star}2$. The fibre part $1-t_0=1-3h$ pays for
$h<\frac13$, for the comparison gap $g$ and for the margins $d_n-h$ and $1-3d_n$ of the transfer; the strip part $\sqrt k\,(a-\rho)$ pays for the enlargement from $\rho$ to $\bar\rho$ and for the excess of the strip map.
The transfer, the ball chart and the dilation are symplectic embeddings and lose no volume beyond these margins.
\end{remark}

